% Standalone section fragment.  Requires amsmath, amssymb, amsthm and
% theorem/lemma/proposition/corollary/example environments in the parent.
\section{Selecting the dual forest by its remaining stopping sets}
\label{sec:forest-choice}

Let $N=(n_{ih})$ be the nonnegative integral occurrence matrix of the
two-handles against a fixed system of one-handle cocores.  Its row set is
$R$.  A subset $E\subseteq R$ labels the edges of a finite undirected
multigraph $D$ with distinguished root $0$.  Rows outside $E$ cannot be
removed by the proposed top-handle cancellations.  Assume that $D$ is
connected, and put $r=|V(D)|-1$.  A successful choice is a spanning tree
$T\subseteq E$ for which the row restriction $N[R\setminus T]$ admits
complete positive raw cancellation.  This section treats that finite
graph--matrix problem; the geometric importer must separately justify the
dual edges and the invariance of the retained source rows.

A nonempty set $S\subseteq R$ is a stopping set if
\[
                 \sum_{i\in S}n_{ih}\ne1\qquad\hbox{for every }h.
\]
The incoming stopping-core theorem proves that complete positive raw
cancellation is equivalent to absence of such a set, and that unit-column
peeling computes the unique largest stopping set.  We use this result as
prior work, including its extension from leaf pivots to all positive pivots.
Replacing each entry by $\min(2,n_{ih})$ preserves the relevant predicates.

\begin{theorem}[Exact forest--stopping transversal criterion]
\label{thm:forest-stopping-transversal}
A spanning tree $T$ is successful if and only if it meets every nonempty
stopping set of the immutable source matrix $N$.  Equivalently, there is
an independent set $F$ in the graphic matroid of $D$ such that
$N[R\setminus F]$ peels completely.  Any such $F$ can be extended to an
arbitrary spanning tree, and every resulting tree is successful.
\end{theorem}
\begin{proof}
A stopping set of a row restriction is exactly a source stopping set
contained in that restriction: its definition sums only its own rows.
The first assertion therefore follows from the imported stopping theorem.
Complete peelability is hereditary under taking row subsets.  Indeed,
any stopping set in a smaller row set would also obstruct the larger
one.  Every independent graph edge set extends to a spanning tree of a
connected graph.  Thus extending $F$ deletes more rows and preserves
peelability.  Conversely a successful tree itself is such an independent
set.
\end{proof}

\subsection{Obstruction-guided branching and independently checkable failure}

Maintain a selected forest $C$.  Contract its graph edges and delete its
matrix rows.  If peeling has empty residual, extend $C$ to a tree and
replay a complete immutable-source peeling schedule.  Otherwise choose
any nonempty stopping set $S$ in the residual.  Every successful extension
must contain an edge of $S$.  An eligible edge that becomes a loop after
contracting $C$ cannot occur in such an extension.  Branch precisely over
the remaining edges of $S$.

\begin{theorem}[Complete constrained search]
\label{thm:forest-branching}
The preceding search is complete for successful dual-forest choices.
Every branch increases $|C|$ by one, so its depth is at most $r$.  If each
chosen witness has at most $w$ eligible nonloop candidates in the current
contraction, the search visits at most
\[
                 1+w+\cdots+w^r
\]
nodes.  A node with no candidate is a proved failure of this restricted
certificate class.  A node with one candidate forces that edge.

For $m=|R|$ rows and $I$ nonzero incidences, an inclusion-minimal stopping
witness can be found with at most $m$ additional sparse peeling calls.
Consequently the complete search takes
$(r+1)\max(1,w)^r\operatorname{poly}(m,I,\Lambda)$ bit operations,
where $\Lambda$ is the full encoded input length.
\end{theorem}
\begin{proof}
Suppose a successful spanning tree extends $C$.  By
Theorem~\ref{thm:forest-stopping-transversal}, it meets $S$ in some edge
$e$.  Since $S$ contains no deleted row, $e\notin C$.  Since a tree is
acyclic, $e$ is a nonloop in $D/C$.  The branch adding $e$ therefore
preserves this successful extension.  Induction establishes completeness.
An empty residual and a verified tree extension give a positive certificate.
If no eligible nonloop remains in $S$, no independent extension can meet
$S$, which proves the negative leaf.  Rank increases at every recursive
step, giving the depth and node bounds.

To shrink a witness, inspect its row labels one at a time.  Remove the
current label tentatively and compute the remaining stopping core.  If
that core is nonempty, replace the witness by it; otherwise keep the
label.  A label that was necessary remains necessary after later
shrinking, since row restrictions cannot create a stopping set.  A label
discarded by an earlier step needs no later visit.  The result is
inclusion-minimal after at most $m$ calls.  This is not an algorithm for
a minimum-cardinality witness, and no such optimality is claimed.
\end{proof}

For a completed negative search, the certificate is a finite proof tree.
An internal node supplies its source stopping set and one child for every
eligible nonloop edge in it; a leaf supplies a stopping set all of whose
eligible edges are loops in the current contraction.  The checker verifies
the stopping sums, the forest along the path, and exact branch coverage.
It need not run the producer's queue, witness minimizer or search policy.
An interrupted search yields no negative certificate.

\begin{proposition}[Disjoint stopping-set rank obstruction]
\label{prop:forest-rank-defect}
At a selected forest $C$, let $S_1,\ldots,S_q$ be pairwise disjoint
nonempty stopping sets in the remaining rows.  Write
$A=E\cap\bigcup_jS_j$.  If
\[
 r_D(C\cup A)-|C|<q,
\]
then no successful tree extends $C$.
\end{proposition}
\begin{proof}
A successful extension must meet each $S_j$.  Disjointness requires at
least $q$ different selected edges of $A$, and they must be independent
over $C$.  The displayed rank defect makes this impossible.
\end{proof}

The source stopping sets, their disjointness and a finite graph rank
computation form a short negative certificate.  A particularly useful
case is a family of mandatory singleton rows whose dual edges form a
cycle.  An ineligible singleton or a loop is already an obstruction.

\subsection{A polynomial regime: each active column meets at most two rows}

Call a column \emph{inert} if none of its entries is one.  Such a column
can never supply a positive raw unit pivot and may be ignored for this
decision, while its one-handle is retained in the topological count.
Assume that every non-inert column has at most two nonzero rows.

Construct a directed graph $P$ on the row set $R$.  A singleton unit
column marks its owner as an initial seed; let $Z$ be the seed set.  For
a column supported on two distinct rows $u,v$, insert $v\to u$ when
$n_{uh}=1$, and insert $u\to v$ when $n_{vh}=1$.  Thus a $(1,1)$ column
gives both directions, a $(1,a)$ column with $a\ge2$ gives only the
direction toward its unit row, and a column with no unit gives none.

\begin{lemma}[Peeling is directed reachability]
\label{lem:two-row-reachability}
After an arbitrary row deletion set $F$, the rows eventually deleted by
column peeling, together with $F$ itself, are exactly the vertices of $P$
reachable from $Z\cup F$.
\end{lemma}
\begin{proof}
A unit leaf comes either from a singleton unit source column, or from a
two-row column whose other owner has already disappeared.  These are
respectively a seed and an enabled directed arc.  Thus each new deletion
is reachable from the stated initial set.

Conversely, a seed can be peeled unless already deleted.  Once $v$ has
disappeared, an arc $v\to u$ exposes a unit column at $u$ unless $u$ has
already disappeared.  That column cannot have been used and removed
while $u$ was still live, because at a leaf pivot it had no other live
owner.  Induction on path length proves the reverse inclusion.
\end{proof}

Let $Q_1,\ldots,Q_p$ be the source strongly connected components of $P$
that contain no seed.  These sets are pairwise disjoint.  Put
$E_j=Q_j\cap E$; graphic loops may be retained in these sets, since they
cannot contribute rank.

\begin{theorem}[Source components are the complete deletion obligations]
\label{thm:source-scc-obligations}
Under the two-row active-column hypothesis, $N[R\setminus F]$ peels
completely if and only if $F$ meets every $Q_j$.  Each $Q_j$ is itself a
source stopping set.
\end{theorem}
\begin{proof}
Contract the strongly connected components of $P$.  Its condensation is
a finite acyclic graph.  Every vertex is reachable from a source vertex.
Reachability from $Z\cup F$ therefore covers the whole graph exactly
when each source component contains a seed or a member of $F$.
Lemma~\ref{lem:two-row-reachability} gives the equivalence.

To check the stopping claim directly, a singleton unit column cannot
meet $Q_j$, since it would give a seed.  If a two-row column had sum one
on $Q_j$, its unit owner would lie in $Q_j$ and its other owner outside.
The resulting incoming arc contradicts the source-component property.
Inert columns never have sum one.  Thus all stopping inequalities hold.
\end{proof}

\begin{theorem}[Polynomial dual-forest selection]
\label{thm:polynomial-dual-forest}
Under the two-row active-column hypothesis, the following are equivalent:
\begin{enumerate}
\item A successful dual spanning tree exists.
\item The disjoint family $(E_j)_{j=1}^p$ has a transversal independent
in the graphic matroid of $D$.
\item For every $J\subseteq\{1,\ldots,p\}$,
\[
                  r_D\left(\bigcup_{j\in J}E_j\right)\ge |J|.
\]
\end{enumerate}
A successful tree and complete raw peeling schedule, or a violated rank
inequality with its pairwise disjoint source stopping sets, can be found
in polynomial encoded time.
\end{theorem}
\begin{proof}
A successful tree meets every $Q_j$.  Choosing one of its edges in each
group gives an independent transversal.  Conversely an independent
transversal extends to a spanning tree.  Theorem~\ref{thm:source-scc-obligations}
and hereditary peelability prove its success.  The rank equivalence is
Rado's classical independent-transversal theorem; an elementary proof is
given below.  Polynomial construction is graphic/partition matroid
intersection.  For completeness, the implemented specialized augmenting
algorithm and its rank obstruction are described next.
\end{proof}

This theorem does not apply matroid intersection to the family of
peelable row sets.  It applies it to graphic independence and the explicit
partition into source-component obligations.  The distinction is necessary.

\begin{example}[Peelable row sets need not be a matroid]
Take rows
\[
 r_0=(0,0,1),\quad r_1=(0,1,1),\quad
 r_2=(1,1,0),\quad r_3=(1,1,1).
\]
The sets $A=\{r_2,r_3\}$ and $B=\{r_0,r_1,r_2\}$ both peel completely.
For $A$, first use the last column on $r_3$; for $B$, use the first column
on $r_2$, then the middle column on $r_1$.  But $A\cup\{r_0\}$ has
column sums $(2,2,2)$ and $A\cup\{r_1\}$ has sums $(2,3,2)$.  Each is
a nonempty stopping set.  Thus $|A|<|B|$ while no member of $B\setminus A$
can be added to $A$, violating matroid exchange.  The middle source
column has three nonzero rows, outside the theorem's hypothesis.
\end{example}

\subsection{The specialized augmenting algorithm and its proof}

Work on $G=\bigcup_j E_j$ and maintain a rainbow forest $F$, meaning a
forest with at most one edge from each group.  Its exchange digraph has
vertices $G$.  The sources are edges $e\notin F$ such that $F+e$ is a
forest.  The targets are edges $e\notin F$ whose group is not represented
in $F$.  For $f\in F$ and $e\notin F$, add
\[
 f\longrightarrow e\quad\hbox{if }F-f+e\hbox{ is a forest},
 \qquad
 e\longrightarrow f\quad\hbox{if }F-f+e\hbox{ is rainbow}.
\]
Choose a shortest source-to-target path and take the symmetric difference
of $F$ with its vertices.  This is the classical matroid-intersection
augmentation specialized to a graph and a partition, not a new general
intersection algorithm.

\begin{lemma}[Shortest augmentation is a larger rainbow forest]
\label{lem:rainbow-augmentation}
The indicated symmetric difference is a rainbow forest with $|F|+1$ edges.
\end{lemma}
\begin{proof}
Write the path as
$e_0,f_1,e_1,\ldots,f_\ell,e_\ell$, where the $e_i$ lie outside $F$.
Before its last vertex no outside edge has an unused group, since it
would already be a target.  Hence $e_{i-1}$ has the group of $f_i$ for
$1\le i\le\ell$, and $e_\ell$ has a new group.  Distinct chosen edges
have distinct groups, so the symmetric difference is rainbow.

Represent graph edges by their vertex-incidence vectors over $\mathbb F_2$.
The vectors of a forest are independent.  The edge $e_0$ is outside the
span of $F$, since it is a source.  For $i\ge1$, the edge $e_i$ is not a
source: otherwise the suffix beginning there would be shorter.  Thus it
has a unique fundamental cycle with $F$, and its expansion in the basis
$F$ has coefficient one at $f_i$.  This expansion has coefficient zero
at $f_j$ for $j<i$, because otherwise $f_j\to e_i$ would shortcut the
chosen path.  Modulo the span of $F\setminus\{f_1,\ldots,f_\ell\}$,
the coefficient matrix of $e_1,\ldots,e_\ell$ on $f_1,\ldots,f_\ell$
is triangular with unit diagonal.  These new edges are independent in
that quotient.  Adding $e_0$, outside the original span, preserves
independence and increases cardinality by one.
\end{proof}

If no target is reachable, let $W$ be the reachable edge set.  Every
color appearing in $W$ is represented by an edge of $F\cap W$: an
outside reachable edge in a new color would be a target, and an outside
edge in an existing color has an exchange arc to that color's chosen
edge.  Thus the number of colors in $W$ is $|F\cap W|$.

Every outside edge $e\notin W$ is spanned by $F\setminus W$.  It is not
a source, and none of the edges of its fundamental cycle can belong to
$F\cap W$, since such an edge would have an exchange arc to $e$.
Consequently
\[
                         r_D(G\setminus W)=|F\setminus W|.
\]
Let $J$ be the groups with no edge in $W$.  Then
$\bigcup_{j\in J}E_j\subseteq G\setminus W$.  If $|F|<p$,
\[
 r_D\left(\bigcup_{j\in J}E_j\right)
 \le |F\setminus W|
 =|F|-|F\cap W|
 <p-|F\cap W|=|J|.
\]
This is the promised short Rado obstruction.  Its independent verifier
only checks the disjoint stopping sets and graph rank, not the exchange
search or its maximality claim.

There are at most $p$ augmentations.  Direct forest oracles, adjacency
search and explicit exchange tests already give polynomial time.  The
reference implementation uses no optimized dynamic-tree data structure:
with $m$ rows and $I$ incidences, a conservative comparison bound is
$O(I\log(I+2)+m(p+1)^3\log(m+2))$, with binary source-label costs added.
This is a polynomial bound for the supplied finite problem.  It is not
a logarithmic hierarchy bound or a diagram-wide complexity theorem.

\subsection{An elementary proof of the rank criterion}

We recall the classical independent-transversal argument to specify its
use precisely.  Let $M$ be a finite matroid and let $E_1,\ldots,E_p$ be
disjoint sets.  Suppose that $r_M(\bigcup_{j\in J}E_j)\ge |J|$ for all
$J$.  Induct on $p$.  The case $p=1$ chooses a nonloop.

If there is a nonempty proper tight set of labels $J$, with
$r_M(\bigcup_{j\in J}E_j)=|J|$, induction chooses a transversal $B$
for these labels.  It is a basis of their union.  In $M/B$, every union
of remaining groups indexed by $K$ has rank
\[
 r_{M/B}\left(\bigcup_{k\in K}E_k\right)
 =r_M\left(\bigcup_{j\in J\cup K}E_j\right)-|J|
 \ge |K|.
\]
Induction supplies the rest of the transversal independently over $B$.
If there is no nonempty proper tight set, choose a nonloop $e\in E_p$.
Contracting $e$ reduces the rank of any union of the other groups by at
most one.  Their original ranks are at least their number of groups
plus one, so the required inequalities persist.  Induction in $M/e$
finishes the proof.  Necessity follows by restricting any independent
transversal to a subset of its labels.

The rank criterion belongs to Rado (1942); the augmenting framework to
Edmonds' matroid-intersection work.  The new application is the
source-component characterization of the raw two/three-handle forest
choice, with its immutable geometric incidence interface.

\subsection{A geometric example where the forest matters}

\begin{example}[Two topologically valid trees, different raw outcomes]
\label{ex:geometric-forest-choice}
Let $H$ be a solid torus with a fixed meridian cocore disc $D$.  Choose
two disjoint parallel primitive curves $a,b$ on $\partial H$.  Arrange
$a$ to meet $\partial D$ once.  In the annulus obtained by cutting along
$a$, introduce a small finger in $b$ across the transverse arc of
$\partial D$, and then back.  The finger is an isotopy supported away
from $a$ and creates two extra transverse intersections of opposite
sign.  Thus $b$ is still embedded, primitive, parallel to and disjoint
from $a$, but the actual unsigned incidences are
\[
                              N=\begin{pmatrix}1\\3\end{pmatrix}.
\]
Attach two-handles along $a$ and $b$.  Compression of the torus along
the two parallel primitive curves yields two sphere boundary components.
Fill the sphere corresponding to the annulus between the curves with
one three-handle.  The resulting manifold is a ball: after cancelling
the primitive first two-handle, the second is attached trivially to a
ball boundary and gives a spherical collar, whose selected sphere is
then filled.

The dual graph has one three-handle vertex, one unfilled boundary root,
and two parallel edges $a,b$.  Both singleton edges are valid spanning
trees.  Deleting the $b$ row leaves $(1)$ and gives full raw cancellation.
Deleting the $a$ row leaves $(3)$, a singleton stopping set.  Both top
cancellations preserve the ball, but only one exposes a certificate in
the fixed raw one-handle system.  The even sphere-incidence parity is
respected, since $1+3$ is even.  The corresponding capped matrix is
$(1,2)^{\mathsf T}$.
\end{example}

This example is a topological existence construction, not a native
triangulated fixture produced by the reference graph module.  The
source-authenticated native tests must supply their own geometric data.

\subsection{A controlled exact separation from tree enumeration}

Take a chain of $k$ dual links, each replaced by two parallel edges
$a_i,b_i$.  Give these two rows the sole nonzero incidences
$n_{a_i,i}=1$ and $n_{b_i,i}=2$ (or actual count three for the second).
The graph has exactly $2^k$ spanning trees, choosing one edge in every
parallel pair.  Every $\{b_i\}$ is a mandatory singleton stopping set,
so the sole successful tree is $\{b_1,\ldots,b_k\}$.

An include-before-exclude tree enumeration with $a_i$ first tests exactly
$2^k$ complete trees before this last success.  The stopping algorithm
uses $k$ forced contractions; the polynomial algorithm has $k$ one-row
source-component obligations.  This proves an exponential separation
from that specified enumeration policy on supplied finite graph/matrix
instances.  It is not a lower bound on arbitrary forest-selection
algorithms, and it does not claim an exponential end-to-end recognition
improvement on knot diagrams.
